\section{Quadratic Voting for Research Funding}

\label{sec:qv}
% ══════════════════════════════════════════════════════════════════════════════

\subsection{Mechanism Design}

In each funding round, the DAO allocates budget $B$ across $N$ proposals.
Voter $i$ receives voice credits $v_i$ proportional to \ZOO{} holdings
(capped). Voter $i$ allocates $k_{ij}$ votes to proposal $j$, paying
$k_{ij}^2$ credits:
\begin{equation}
  \sum_{j=1}^{N} k_{ij}^2 \leq v_i \quad \forall i
\end{equation}

Total votes for proposal $j$:
\begin{equation}
  V_j = \sum_{i=1}^{M} k_{ij}
\end{equation}

Budget allocation proportional to squared total votes:
\begin{equation}
  B_j = B \cdot \frac{\max(V_j, 0)^2}{\sum_{j'} \max(V_{j'}, 0)^2}
\end{equation}

Negative votes are permitted to signal opposition.

\subsection{Sybil Resistance}

We implement three defenses: ZooRep gating ($\geq$50 ZooRep to vote),
pairwise coordination detection \citep{buterin2019flexible}, and
connection-weighted QV that scales credits by social graph uniqueness.

\subsection{Formal Properties}

\begin{theorem}[Efficiency of QV for Research Allocation]
Under quasi-linear utility over research outcomes and sufficient voice
credits, QV funding allocation converges to the utilitarian
welfare-maximizing allocation as voter count increases.
\end{theorem}

\begin{proof}
Under quasi-linear utility $u_i(B_j) = \theta_{ij} \cdot B_j -
k_{ij}^2$ where $\theta_{ij}$ is voter $i$'s marginal valuation, the
first-order condition gives $k_{ij}^* = \theta_{ij} / 2$. The aggregate
vote $V_j = \sum_i \theta_{ij} / 2$ is proportional to total social
valuation. The QF rule ensures funding proportional to the square of
aggregate valuation, approximating the utilitarian optimum
\citep{lalley2018quadratic}.
\end{proof}

\subsection{Comparison with Alternatives}

\begin{table}[H]
\centering
\caption{Comparison of voting mechanisms for research funding.}
\label{tab:voting}
\begin{tabular}{p{3cm}cccc}
\toprule
\textbf{Property} & \textbf{1T1V} & \textbf{QV} &
  \textbf{Conviction} & \textbf{Expert panel} \\
\midrule
Plutocracy resistance & Low & High & Medium & N/A \\
Preference intensity & No & Yes & Partial & Yes \\
Sybil resistance & High & Medium & High & High \\
Participation cost & Low & Medium & Low & High \\
Minority voice & Low & High & Medium & Low \\
Scalability & High & High & High & Low \\
\bottomrule
\end{tabular}
\end{table}


% ══════════════════════════════════════════════════════════════════════════════
